\chapter{Introduction to Game Programming in Python (Pygame)}

\section{Introduction: Real-Life Hook}

Think about your favorite video games:  

\begin{itemize}
    \item Moving a character across the screen  
    \item Collecting points or avoiding obstacles  
    \item Interactive buttons and score displays  
\end{itemize}

Python’s \textbf{Pygame} library allows you to \textbf{create simple 2D games}, helping you combine \textbf{loops, functions, events, and graphics} in a fun way.

\section{Concept: What is Pygame?}

\begin{definition}{Pygame}
Pygame is a Python library that helps you create \textbf{games and multimedia applications} using graphics, sounds, and events.
\end{definition}

\section{Installing Pygame}

Use pip to install:

\begin{verbatim}
pip install pygame
\end{verbatim}

\section{Creating a Basic Game Window}

\begin{verbatim}
import pygame

# Initialize Pygame
pygame.init()

# Create window
window = pygame.display.set_mode((400, 300))
pygame.display.set_caption("My First Game")

running = True
while running:
    for event in pygame.event.get():
        if event.type == pygame.QUIT:
            running = False

pygame.quit()
\end{verbatim}

- \texttt{pygame.init()} initializes all Pygame modules  
- \texttt{set\_mode()} creates the game window  
- Event loop listens for actions like closing the window  

\section{Drawing Shapes and Colors}

\begin{verbatim}
import pygame

pygame.init()
window = pygame.display.set_mode((400, 300))
pygame.display.set_caption("Draw Shapes")

# Colors
white = (255, 255, 255)
blue = (0, 0, 255)

# Fill background
window.fill(white)

# Draw a rectangle
pygame.draw.rect(window, blue, (50, 50, 100, 50))  # x, y, width, height

pygame.display.update()  # Refresh display

# Wait for 3 seconds
pygame.time.wait(3000)
pygame.quit()
\end{verbatim}

- Use RGB tuples for colors  
- \texttt{pygame.draw} allows drawing rectangles, circles, lines, etc.  

\section{Handling Keyboard Events}

\begin{verbatim}
import pygame

pygame.init()
window = pygame.display.set_mode((400, 300))
pygame.display.set_caption("Keyboard Example")

running = True
while running:
    for event in pygame.event.get():
        if event.type == pygame.QUIT:
            running = False
        elif event.type == pygame.KEYDOWN:
            if event.key == pygame.K_UP:
                print("Up arrow pressed")
            elif event.key == pygame.K_DOWN:
                print("Down arrow pressed")

pygame.quit()
\end{verbatim}

- Detect which key the player presses  
- Useful for moving characters or controlling actions  

\section{Mini-Project: Catch the Ball Game}

Concept:

\begin{itemize}
    \item Ball falls from the top of the window  
    \item Player controls a basket at the bottom using left/right keys  
    \item Catch the ball to earn points  
\end{itemize}

Basic structure:

\begin{verbatim}
import pygame
import random

pygame.init()
window = pygame.display.set_mode((400, 500))
pygame.display.set_caption("Catch the Ball")
clock = pygame.time.Clock()

# Colors
white = (255,255,255)
red = (255,0,0)
blue = (0,0,255)

# Basket
basket_width = 80
basket_x = 160
basket_y = 450

# Ball
ball_radius = 20
ball_x = random.randint(ball_radius, 400-ball_radius)
ball_y = 0
ball_speed = 5

score = 0
running = True
while running:
    window.fill(white)

    # Draw basket
    pygame.draw.rect(window, blue, (basket_x, basket_y, basket_width, 20))

    # Draw ball
    pygame.draw.circle(window, red, (ball_x, ball_y), ball_radius)

    # Move ball
    ball_y += ball_speed

    # Check catch
    if (basket_y <= ball_y + ball_radius <= basket_y + 20) and 
       (basket_x <= ball_x <= basket_x + basket_width):
        score += 1
        ball_y = 0
        ball_x = random.randint(ball_radius, 400-ball_radius)

    # Event handling
    for event in pygame.event.get():
        if event.type == pygame.QUIT:
            running = False

    # Move basket
    keys = pygame.key.get_pressed()
    if keys[pygame.K_LEFT] and basket_x > 0:
        basket_x -= 5
    if keys[pygame.K_RIGHT] and basket_x < 320:
        basket_x += 5

    # Display score
    font = pygame.font.Font(None, 36)
    text = font.render(f"Score: {score}", True, (0,0,0))
    window.blit(text, (10,10))

    pygame.display.update()
    clock.tick(30)

pygame.quit()
\end{verbatim}

- Students learn \textbf{graphics, movement, collision detection, and scoring}  

\section{Step-by-Step Diagram: Pygame Game Loop}

\begin{center}
\begin{tikzpicture}[node distance=1.5cm]
\node (start) [draw, ellipse, fill=green!20] {Start Game};
\node (event) [draw, rectangle, fill=yellow!20, below of=start] {Check Events (Keyboard/Mouse)};
\node (update) [draw, rectangle, fill=orange!20, below of=event] {Update Game State (Move Ball/Basket)};
\node (draw) [draw, rectangle, fill=yellow!20, below of=update] {Draw Graphics on Screen};
\node (loop) [draw, diamond, aspect=2, fill=orange!20, below of=draw, yshift=-0.5cm] {Loop Until Game Ends};
\node (end) [draw, ellipse, fill=green!20, below of=loop, yshift=-1cm] {End Game};
\draw[->] (start) -- (event);
\draw[->] (event) -- (update);
\draw[->] (update) -- (draw);
\draw[->] (draw) -- (loop);
\draw[->] (loop.south) -- (end.north);
\draw[->] (loop.east) -- ++(2,0) |- (event.east);
\end{tikzpicture}
\end{center}

\section{Key Concepts}

\begin{itemize}
    \item Pygame allows \textbf{2D graphics and multimedia} in Python  
    \item Game loop: \textbf{event handling $\rightarrow$ update state $\rightarrow$ draw graphics}  
    \item Keyboard and mouse events control game objects  
    \item Collision detection is essential for interaction  
    \item Scoring and game state updates make games interactive  
\end{itemize}

\section{End-of-Chapter Exercises}

\begin{exercise}{Very Short Questions}
\begin{enumerate}
    \item What is Pygame?  
    \item What is a game loop?  
    \item Name two Pygame functions for drawing shapes.  
\end{enumerate}
\end{exercise}

\begin{exercise}{Short Answer Questions}
\begin{enumerate}
    \item Write code to create a Pygame window with a moving rectangle.  
    \item Explain how \texttt{pygame.key.get\_pressed()} works.  
    \item What is the purpose of \texttt{clock.tick()} in the game loop?  
\end{enumerate}
\end{exercise}

\begin{exercise}{Activity-Based Questions}
\begin{enumerate}
    \item Modify the Catch the Ball game to include multiple balls falling at different speeds.  
    \item Add a timer to the game that ends after 60 seconds and displays the final score.  
    \item Create a simple maze game where a player moves a square using arrow keys and reaches a goal.  
\end{enumerate}
\end{exercise}

\section{Teacher Notes}

\marginnote{
\textbf{Learning Objectives:}  
\begin{itemize}
    \item Students can create windows, draw shapes, and handle keyboard events  
    \item Students understand the game loop and event-driven programming  
    \item Students can build simple interactive games with scoring  
\end{itemize}

\textbf{Teaching Tips:}  
\begin{itemize}
    \item Start with basic window and shapes before adding movement  
    \item Encourage students to experiment with colors, speed, and sizes  
    \item Show step-by-step collision detection and scoring logic  
\end{itemize}

\textbf{Assessment:}  
\begin{itemize}
    \item Check correctness of game loop and event handling  
    \item Observe interactive elements, movement, and scoring  
\end{itemize}
}